%% sections/08-troubleshooting.tex
%% Troubleshooting Guide
%% IEC/IEEE 82079-1:2019 — Reference Information (§5)
%% ─────────────────────────────────────────────────────

\section{Troubleshooting}
\label{sec:troubleshooting}

%%
%% WRITING GUIDE — TROUBLESHOOTING:
%%
%%   Use the \begin{troubletable} environment for the main table.
%%   Each row: \troubrow{Problem}{Probable Cause}{Corrective Action}
%%
%%   Keep corrective actions specific and actionable.
%%   Use \code{}, \acsparam{}, \cref{} inside cells.
%%   Add rows as new issues are discovered in the field.
%%

\begin{troubletable}

  \troubrow{%
    Axis does not enable; MFLAGS bit 0 not set
  }{%
    E-Stop active, or interlock circuit open
  }{%
    Confirm E-Stop is released.  Check interlock status in
    \menu{Axes > I/O Status}.
  }

  \troubrow{%
    Force reading shows large offset ($>$\SI{0.5}{\newton}) at zero load
  }{%
    Sensor not zeroed, or temperature drift
  }{%
    Re-run Phase~1, Step~2 (\cref{sec:proc:phase1}).
    Check sensor temperature stabilisation time ($\geq$15~min warm-up).
  }

  \troubrow{%
    High noise on force signal ($>$\SI{0.5}{\newton} RMS)
  }{%
    Ground loop on sensor cable shield, or nearby EMI source
  }{%
    Verify cable shield is grounded at conditioner end ONLY.
    Route cable away from motor drive cables.
    Check AIN differential mode in MMI: \menu{Axes > Analogue Inputs}.
  }

  \troubrow{%
    Force loop unstable — oscillates continuously
  }{%
    \acsparam{KP} too high, or derivative gain \acsparam{KD} too low
  }{%
    Reduce \acsparam{KP} by 50\,\%.  Set \acsparam{KI} = 0 temporarily.
    Retune from baseline (\cref{tab:pid-defaults}).
  }

  \troubrow{%
    Force undershoots setpoint; very slow convergence
  }{%
    \acsparam{KI} too low, or current limit \acsparam{DCURR} too low
  }{%
    Increase \acsparam{KI} by 0.01 steps.
    Verify \acsparam{DCURR} is set to the motor's rated continuous
    current (see motor datasheet).
  }

  \troubrow{%
    ACS axis fault during force control run (FAULTS bit set)
  }{%
    Position error fault triggered (PEF), or current limit exceeded
  }{%
    Clear fault: \code{CLEAR} in MMI terminal.
    Check \acsparam{XVEL} is set to a safe velocity limit.
    Check workpiece compliance — very stiff contact surfaces can cause
    transient current spikes.
  }

  \troubrow{%
    Python script cannot connect to ACS controller
  }{%
    Wrong IP address, or firewall blocking port 701
  }{%
    Run \code{ping \textit{<controller-ip>}}.  Confirm port 701 is not
    blocked by Windows Firewall.  Update \code{ACS\_IP} in the script.
    See \cref{sec:setup:network}.
  }

  \troubrow{%
    CSV force log shows only zeros
  }{%
    \code{acspy} library not connected, or wrong channel index
  }{%
    Check the \code{channel} parameter in the Python script matches the
    physical AIN port used (0-indexed: AIN~1 = index~0).
  }

  \troubrow{%
    Force drifts upward over time during sustained contact
  }{%
    Thermal zero drift of load cell, or material creep in workpiece
  }{%
    Re-zero sensor every 30~min.
    Record ambient temperature alongside force data.
    Apply thermal compensation factor from sensor datasheet.
  }

  \troubrow{%
    ACS MMI shows ``Configuration mismatch'' after loading \filepath{.cfg}
  }{%
    Firmware version mismatch between config file and controller
  }{%
    Confirm controller firmware is $\geq$ 2.85.
    Contact PBA Systems support with controller serial number to obtain
    the correct config version.
  }

\end{troubletable}

\begin{tipbox}
  Record all field-discovered issues and their resolutions in the FORGE
  dead-end registry under project \texttt{PBA-GANTRY-FC}.
  This keeps the troubleshooting table growing with every deployment.
\end{tipbox}
